% Source: convolutional_networks.tex
\begin{notechapter}
\chapter{Convolution in CNNs}\label{ch:convolution-in-cnns}

\section{Introduction}

Convolutional Neural Networks (CNNs) are a cornerstone of modern deep 
learning, particularly for image processing tasks. This document 
explains how convolution works in CNNs, with a focus on RGB images, 
and provides intuitive examples to solidify understanding.

\section{Standard Convolution}

\subsection{Input and Kernel Structure}

In CNNs, an RGB image is represented as a 3D tensor of size 
\(H \times W \times C\), where:
\begin{itemize}
    \item \(H\) is the height of the image,
    \item \(W\) is the width of the image,
    \item \(C\) is the number of channels (e.g., \(C = 3\) for RGB).
\end{itemize}

A convolutional kernel (or filter) is also a 3D tensor of size 
\(k \times k \times C\), where:
\begin{itemize}
    \item \(k \times k\) is the spatial size of the kernel,
    \item \(C\) matches the number of input channels.
\end{itemize}

For example, a \(3 \times 3\) kernel applied to an RGB image has a 
shape of \(3 \times 3 \times 3\).

\subsection{Convolution Operation}

At each spatial location \((x, y)\) in the input, the kernel is 
applied by performing element-wise multiplication between the kernel 
and the corresponding patch of the input feature maps. The results 
are summed across all \(C\) input channels to produce a single scalar 
value:

\[
\text{Output}(x, y) = \sum_{i=1}^{k} \sum_{j=1}^{k} \sum_{c=1}^{C} 
\text{Kernel}(i, j, c) \cdot \text{Input}(x+i, y+j, c)
\]

This operation is repeated for every spatial location, resulting in 
a single 2D output feature map.

\subsection{Example: Convolution on an RGB Image}

Consider a \(3 \times 3\) RGB image with 3 channels. Let's focus on 
the entire image:

\[
\text{Input}_{\text{Red}} = \begin{bmatrix}
1 & 2 & 3 \\
4 & 5 & 6 \\
7 & 8 & 9
\end{bmatrix}, \quad
\text{Input}_{\text{Green}} = \begin{bmatrix}
9 & 8 & 7 \\
6 & 5 & 4 \\
3 & 2 & 1
\end{bmatrix}, \quad
\text{Input}_{\text{Blue}} = \begin{bmatrix}
2 & 4 & 6 \\
8 & 1 & 3 \\
5 & 7 & 9
\end{bmatrix}
\]

Let the kernel be:

\[
\text{Kernel}_{\text{Red}} = \begin{bmatrix}
1 & 0 & -1 \\
1 & 0 & -1 \\
1 & 0 & -1
\end{bmatrix}, \quad
\text{Kernel}_{\text{Green}} = \begin{bmatrix}
0 & 1 & 0 \\
-1 & 0 & 1 \\
0 & -1 & 0
\end{bmatrix}, \quad
\text{Kernel}_{\text{Blue}} = \begin{bmatrix}
-1 & 0 & 1 \\
0 & 0 & 0 \\
1 & 0 & -1
\end{bmatrix}
\]

The output at position \((1, 1)\) is computed as:

\[
\text{Output}(1, 1) = \sum_{c=1}^{3} \left( \sum_{i=1}^{3} \sum_{j=1}^{3} 
\text{Kernel}(i, j, c) \cdot \text{Input}(i, j, c) \right)
\]

For the Red channel:

\[
(1 \cdot 1) + (0 \cdot 2) + (-1 \cdot 3) + (1 \cdot 4) + (0 \cdot 5) 
+ (-1 \cdot 6) + (1 \cdot 7) + (0 \cdot 8) + (-1 \cdot 9) = -6
\]

Similarly, compute contributions from Green and Blue channels, then 
sum them to get the final output value.

\subsection{Why Combine Contributions?}

Combining contributions from all input channels allows the network 
to learn \textbf{cross-channel relationships}. For example, in an RGB 
image, the kernel can detect patterns involving interactions between 
the Red, Green, and Blue channels. This is crucial for capturing 
complex features like edges, textures, and shapes.

\section{Depthwise Separable Convolution}

Depthwise separable convolution is a more efficient alternative to 
standard convolution, used in architectures like MobileNet. It 
decomposes the standard convolution into two steps: \textbf{depthwise 
convolution} and \textbf{pointwise convolution}.

\subsection{Depthwise Convolution}

In depthwise convolution, each input channel is convolved with its 
own separate 2D kernel (\(k \times k \times 1\)). This means:
\begin{itemize}
    \item If the input has \(C\) channels, there are \(C\) separate 
    2D kernels, each of size \(k \times k\).
    \item Each kernel is applied independently to its corresponding 
    input channel.
    \item The result is \(C\) intermediate output feature maps, each 
    of which is the result of convolving one input channel with its 
    corresponding kernel.
\end{itemize}

\subsubsection{Key Points:}
\begin{itemize}
    \item \textbf{No Cross-Channel Interaction}: Each output feature 
    map depends only on its corresponding input channel. There is no 
    mixing of information across channels at this stage.
    \item \textbf{Output Shape}: If the input has shape \(H \times W 
    \times C\), the output of the depthwise convolution will have 
    shape \(H' \times W' \times C\), where \(H'\) and \(W'\) depend 
    on the padding and stride.
\end{itemize}

\subsubsection{Mathematical Representation:}
For each channel \(c\), the output feature map is computed as:
\[
\text{DepthwiseOutput}(x, y, c) = \sum_{i=1}^{k} \sum_{j=1}^{k} 
\text{Kernel}_{c}(i, j) \cdot \text{Input}(x+i, y+j, c)
\]
Here, \(\text{Kernel}_{c}\) is the 2D kernel for channel \(c\).

\subsection{Pointwise Convolution}

After the depthwise convolution, \textbf{pointwise convolution} is 
applied to combine the intermediate output feature maps across 
channels. This is done using \(1 \times 1\) kernels:
\begin{itemize}
    \item Each \(1 \times 1\) kernel has \(C\) weights, one for each 
    intermediate output feature map.
    \item The pointwise convolution produces a new set of output 
    feature maps, where each output feature map is a linear 
    combination of the intermediate output feature maps.
\end{itemize}

\subsubsection{Key Points:}
\begin{itemize}
    \item \textbf{Cross-Channel Interaction}: The pointwise 
    convolution allows the network to learn relationships between 
    channels.
    \item \textbf{Output Shape}: If \(D\) \(1 \times 1\) kernels are 
    used, the output will have shape \(H' \times W' \times D\).
\end{itemize}

\subsubsection{Mathematical Representation:}
For each output feature map \(d\), the output is computed as:
\[
\text{Output}(x, y, d) = \sum_{c=1}^{C} \text{Kernel}_{1 \times 1}(c, d) 
\cdot \text{DepthwiseOutput}(x, y, c)
\]
Here, \(\text{Kernel}_{1 \times 1}(c, d)\) is the weight for channel 
\(c\) in the \(1 \times 1\) kernel for output feature map \(d\).

\subsection{Example: Depthwise Separable Convolution on an RGB Image}

Consider the same \(3 \times 3\) RGB image with 3 channels. Let's 
apply depthwise separable convolution:

\subsubsection{Depthwise Convolution:}
\begin{itemize}
    \item Use 3 separate \(3 \times 3\) kernels, one for each channel.
    \item For each channel, compute the intermediate output feature map.
\end{itemize}

For the Red channel:

\[
\text{DepthwiseOutput}_{\text{Red}} = \sum_{i=1}^{3} \sum_{j=1}^{3} 
\text{Kernel}_{\text{Red}}(i, j) \cdot \text{Input}_{\text{Red}}(i, j)
\]

Substituting the values:

\[
(1 \cdot 1) + (0 \cdot 2) + (-1 \cdot 3) + (1 \cdot 4) + (0 \cdot 5) 
+ (-1 \cdot 6) + (1 \cdot 7) + (0 \cdot 8) + (-1 \cdot 9) = -6
\]

Similarly, compute intermediate outputs for Green and Blue channels.

\subsubsection{Pointwise Convolution:}
\begin{itemize}
    \item Use \(D\) \(1 \times 1\) kernels to combine the intermediate outputs.
    \item For example, if \(D = 2\), the final output will have 2 feature maps.
\end{itemize}

For the first output feature map (\(d = 1\)):

\[
\text{Output}(1, 1, 1) = \sum_{c=1}^{3} \text{Kernel}_{1 \times 1}(c, 1) 
\cdot \text{DepthwiseOutput}(1, 1, c)
\]

For the second output feature map (\(d = 2\)):

\[
\text{Output}(1, 1, 2) = \sum_{c=1}^{3} \text{Kernel}_{1 \times 1}(c, 2) 
\cdot \text{DepthwiseOutput}(1, 1, c)
\]

\section{Summary}

\begin{itemize}
    \item In \textbf{standard convolution}, contributions from all 
    input channels are combined to produce a single output value.
    \item In \textbf{depthwise separable convolution}:
    \begin{itemize}
        \item \textbf{Depthwise Convolution}: No linear combination 
        across channels. Each output feature map depends only on its 
        corresponding input channel.
        \item \textbf{Pointwise Convolution}: Linear combination 
        across channels. Each output feature map is a weighted sum 
        of all intermediate output feature maps.
    \end{itemize}
\end{itemize}

\section{Visual Analogy}

Imagine the kernel as a small flashlight that scans the image. For 
standard convolution, the flashlight has 3 colored bulbs (Red, Green, 
Blue) that shine together, blending their light to produce a single 
output. For depthwise separable convolution, each bulb shines 
independently first, and then their lights are mixed afterward.
\end{notechapter}
